\documentclass[11pt,letterpaper]{article}
\usepackage[margin=1in]{geometry}
\usepackage[T1]{fontenc}
\usepackage{lmodern,microtype}
\usepackage{amsmath,amssymb,booktabs,tabularx}
\usepackage{xcolor,enumitem,titlesec,fancyhdr}
\usepackage[font=small,labelfont=bf]{caption}
\usepackage[numbers,sort&compress]{natbib}
\usepackage[hidelinks]{hyperref}
\usepackage{xurl}
\definecolor{ink}{HTML}{243746}
\definecolor{muted}{HTML}{646B72}
\titleformat{\section}{\large\bfseries\color{ink}}{\thesection}{0.7em}{}
\titleformat{\subsection}{\normalsize\bfseries}{\thesubsection}{0.7em}{}
\setlength{\parindent}{1.1em}
\setlength{\parskip}{0.15em}
\setlength{\headheight}{14pt}
\setlist{nosep,leftmargin=1.5em}
\renewcommand{\arraystretch}{1.16}
\pagestyle{fancy}
\fancyhf{}
\fancyhead[L]{\small\color{muted}Action semantics for instructional video retrieval}
\fancyfoot[C]{\small\thepage}
\renewcommand{\headrulewidth}{0.2pt}
\newcommand{\pending}[1]{\par\noindent\textcolor{muted}{\textit{[To be completed: #1]}}\par\vspace{0.6em}}
% Source: project1_outputs/experiment-suite-2026-10-06
% Summary SHA-256: 537ad831820fbbd02d9ca1f0533fd0b212a0d408db317a795c7ec443b3be4292
\newcommand{\VideoCount}{250}
\newcommand{\ClipCount}{1,663}
\newcommand{\QueryCount}{462}
\newcommand{\CandidateCount}{582}
\newcommand{\LexicalHitOne}{63.2}
\newcommand{\StructuredHitOne}{26.4}
\newcommand{\HybridHitOne}{50.6}
\newcommand{\InventoryRows}{
Parent videos & 250 \\
Raw clip annotations & 1,703 \\
Valid annotations & 1,700 \\
Rejected intervals & 3 \\
Canonical clips & 1,663 \\
Extracted action records & 8,899 \\
Clips with at least one extracted action & 1,296 (77.9\%) \\
Action records with an object & 64.0\% \\
Action records with a directly attached tool & 7.4\% \\
}
\newcommand{\GlobalRows}{
Lexical TF--IDF & 63.20 & 85.50 & 93.29 & 0.749 \\
Structured & 26.43 & 36.31 & 47.05 & 0.335 \\
Hybrid (50/50) & 50.65 & 61.48 & 67.96 & 0.577 \\
}
\newcommand{\WithinRows}{
Lexical TF--IDF & 72.73 & 92.86 & 98.70 & 0.834 \\
Structured & 50.42 & 68.46 & 73.59 & 0.597 \\
Hybrid (50/50) & 72.73 & 92.50 & 99.78 & 0.833 \\
}
\newcommand{\AblationRows}{
Cleaned baseline & 63.20 & 26.43 & 50.65 \\
Aliases only & 63.20 & 26.43 & 50.65 \\
Primary inventory only & 61.90 & 26.36 & 50.65 \\
Both changes & 61.47 & 26.36 & 50.65 \\
}
\newcommand{\DeltaRows}{
Aliases only & 0.00 & -0.62 & 0.65 \\
Primary inventory only & -1.30 & -3.04 & 0.45 \\
Both changes & -1.73 & -3.66 & 0.21 \\
}
\newcommand{\PairwiseRows}{
Metadata contrast & Lexical & 2,395 & 99.83 & [99.66, 99.96] \\
Metadata contrast & Structured & 2,395 & 95.91 & [95.08, 96.72] \\
Metadata contrast & Hybrid & 2,395 & 99.75 & [99.54, 99.92] \\
Hard negatives & Lexical & 1,309 & 96.72 & [95.36, 97.92] \\
Hard negatives & Structured & 1,309 & 50.76 & [48.24, 53.31] \\
Hard negatives & Hybrid & 1,309 & 72.80 & [70.94, 74.73] \\
}
\newcommand{\DiagnosticRows}{
Lexical/structured disagreement & 172 & 858 \\
Structured ties & 144 & 413 \\
Zero structured evidence & 46 & 369 \\
Query parse failure & 21 & 0 \\
Hybrid recovery relative to lexical & 2 & 13 \\
Hybrid regression relative to lexical & 4 & 326 \\
}

\title{\vspace{-1.8em}\textbf{Action Semantics and Terminology Normalization\\for Instructional Video Retrieval}}
\author{\textit{Author and affiliation to be added}}
\date{Working manuscript\quad October 15, 2026}
\hypersetup{pdftitle={Action Semantics and Terminology Normalization for Instructional Video Retrieval},
  pdfsubject={Development experiments and evaluation protocol},pdfauthor={}}

\begin{document}
\maketitle
\thispagestyle{empty}
\begin{abstract}
Instructional video search requires more than finding clips about the same
topic. A useful result should show the requested action on the appropriate
object. This study compares lexical TF--IDF matching, structured
action--object--tool matching, and an equal-weight hybrid using human-written
annotations from \VideoCount{} videos. Validation and deduplication produce
\ClipCount{} canonical clips. A field-held-out development benchmark evaluates
\QueryCount{} clip-name queries against \CandidateCount{} candidates whose
representations exclude clip names and parent-video text. Lexical matching
achieves \LexicalHitOne\% Hit@1, compared with \StructuredHitOne\% for
structured matching and \HybridHitOne\% for the hybrid. Two constructed
pairwise datasets also favor lexical matching. Inventory parsing fixes remove
identified formatting defects, while a small terminology lexicon yields no
aggregate Hit@1 improvement. These results establish a working evaluation
pipeline and identify weaknesses in the current structured representation.
Independent human preference evaluation and final held-out validation remain
incomplete.
\end{abstract}

\section{Introduction}
A request such as ``remove the old faucet'' is specific about both an action
and an object. A clip about installing a faucet may share much of the same
vocabulary while answering a different question. This motivates representing
what an instruction asks someone to do, rather than relying only on shared
words.

The central question is whether an explicit action--object--tool
representation improves matching between a project step and an instructional
clip. The present implementation compares this representation with a lexical
baseline and a fixed hybrid. It also examines a practical data issue:
different names can refer to the same object or tool. For example, an
annotation may use ``plasterboard'' where a query uses ``drywall.'' Human
authorship establishes the source of the annotations, but does not guarantee
consistent terminology or complete fields.

This manuscript reports the completed development experiments. It separates
three questions: whether the pipeline runs correctly, whether a scorer
recovers a target defined by the benchmark, and whether the scorer agrees
with independent human judgments. The current evidence answers the first two.
It does not yet answer the third.

\section{Related work}
\pending{A focused review of instructional video retrieval, lexical and
semantic retrieval baselines, and terminology normalization. Establish the
relationship to prior methods before making a claim of novelty.}

\section{Data and preparation}
\subsection{Source annotations and canonical clips}
The supplied indexed-video export contains authentic human-generated
annotations, as confirmed by the data owner. Each video contains timestamped
clip records with a name and optional description, goal, tools, and supplies.
The study searches these textual records; it does not analyze video frames
or audio.

The preparation stage rejects non-positive time intervals and groups repeated
annotations with the same source video and timestamps. It merges useful field
variants while preserving the original records. Canonical identifiers depend
on the source video and interval, not the wording of an annotation.
Table~\ref{tab:data} summarizes the resulting index.

\begin{table}[htbp]
\centering\small
\caption{Data inventory and extraction coverage. Coverage measures the presence
of extracted information, not its correctness.}\label{tab:data}
\begin{tabularx}{\linewidth}{Xr}\toprule
Measure & Value\\\midrule
\InventoryRows
\bottomrule\end{tabularx}
\end{table}

Of the \ClipCount{} canonical clips, 1,081 lack descriptions and 1,188 lack
goals. Consequently, the retrieval benchmark uses a smaller candidate set
with narrative evidence. Its results should not be interpreted as accuracy
over every clip in the full index.

\subsection{Annotation provenance and evaluation labels}
Source authorship and evaluation-label construction are distinct. The source
descriptions are human-written. The retrieval target is assigned by the
relationship between fields belonging to the same clip. In the two development
pairwise datasets, a construction rule supplies the winning clip. Neither
procedure establishes that a person compared the candidates and preferred
that winner.

\subsection{Collection and annotation procedure}
\pending{Document video sampling, annotator instructions, authoring and review
procedures, and any restrictions on releasing the data. These details have
not been established from the local export.}

\section{Methods}
\subsection{Lexical baseline}
The lexical method represents text using TF--IDF features for individual words
and adjacent word pairs. The vectors are normalized to unit length, so their
dot product gives cosine similarity~\cite{sklearn}. The retrieval benchmark
fits the vocabulary on candidate fields only and uses raw term frequency.
The existing search and pairwise-evaluation path uses sublinear term
frequency, $1+\log(\mathrm{tf})$, for positive counts. These configurations
are retained as implemented rather than changed after observing results.

The candidate fields also differ by evaluation. The field-held-out benchmark
excludes clip names and parent-video text. The frozen pairwise evaluator uses
the existing search representation, which includes those fields and weights
the clip title twice. Direct comparison of percentages across these two
experiments would therefore be inappropriate.

\subsection{Structured matching}
The parser uses spaCy tokenization, part-of-speech tagging, lemmatization, and
dependency relationships~\cite{spacy}. It extracts action records with an
action lemma, object terms, directly attached tool terms, and negation.
Tool and supply inventories provide separate record-level context. An
inventory item is not assumed to be grammatically attached to every action.

An exact action-lemma match receives a score of one. Otherwise, shared
VerbNet classes and FrameNet frames provide weaker matching evidence.
VerbNet groups verbs using syntactic and semantic patterns~\cite{verbnet};
FrameNet supplies frame-semantic descriptions~\cite{baker1998framenet}.
For sets of labels $X$ and $Y$, let
\begin{equation}
 J(X,Y)=\frac{|X\cap Y|}{|X\cup Y|},
\end{equation}
with zero assigned when there is no shared evidence. If $V$ denotes VerbNet
classes and $F$ denotes FrameNet frames, action similarity is
\begin{equation}
 A(q,c)=
 \begin{cases}
  1, & \text{if the action lemmas match},\\
  \max\{0.80J(V_q,V_c),\;0.70J(F_q,F_c)\}, & \text{otherwise}.
 \end{cases}
\end{equation}
Object matching measures the fraction of query object terms found in the
candidate. Tool and supply comparisons use the same directional coverage
principle. Their strongest applicable match supplies a context score $C$.

For one query action and one candidate action, the score has the form
\begin{equation}\label{eq:structured}
 S(q,c)=\gamma A(q,c)
       \bigl(w_A A(q,c)+w_O O(q,c)+w_C C(q,c)\bigr),
\end{equation}
where the base weights are $(0.55,0.35,0.10)$. Weights for absent query object
or context evidence are removed, and the remaining weights are renormalized.
The confidence multiplier is $\gamma=0.90+0.10\min(p_q,p_c)$, where the $p$
values are extraction confidences. A negation mismatch forces the score to
zero. The strongest candidate action is selected for each query action, and
the resulting scores are averaged. A learned exploratory action taxonomy is
recorded for diagnostics but does not affect ranking.

\subsection{Hybrid matching}
The fixed hybrid uses
\begin{equation}
 H(q,c)=0.5L(q,c)+0.5S(q,c),
\end{equation}
where $L$ is the lexical score. The search and pairwise paths fall back to the
lexical score if no query action can be parsed. The retrieval benchmark instead
excludes such queries while retaining their clips as possible distractors.
No hybrid-weight sweep is used to select the reported result.

\subsection{Inventory cleanup and terminology normalization}
The audit found that comma splitting sometimes broke valid parenthesized
inventory lists. Parsed names also retained placeholders such as ``Unknown''
and ``Unspecified.'' The revised parser respects grouped lists, preserves
usage notes, and removes these placeholders from parsed names. Original
annotations remain available for inspection. On the complete canonical
index, mismatched-parenthesis flags fell from 111 to zero, and placeholder-name
flags fell from 91 to zero. These counts describe audit flags, not a manually
estimated error rate.

The active lexicon contains nine directional rewrites covering six preferred
terms. Examples include ``gypsum board'' and ``plasterboard'' to ``drywall,''
``measuring tape'' to ``tape measure,'' and ``mitre saw'' to ``miter saw.''
The same normalization is applied to queries and candidate representations.
Attributes and original text are retained. Ambiguous terms such as ``tap,''
``mud,'' and ``Sheetrock'' have no global rewrite rule. The local annotations,
for example, use Sheetrock in the name of a joint-compound product.

A second, independent option excludes explicitly listed substitutes from
inventory evidence. It does not remove the original alternatives from the
metadata. This distinction is needed because a circular saw may be offered as
an alternative to a miter saw without being the same tool.

\section{Experimental design}
\subsection{Field-held-out retrieval}
The query is a clip name, and the known target is the same canonical clip.
Candidates are represented by descriptions, goals, tools, and supplies.
Candidates require a description or goal. Queries are excluded when their
normalized names are ambiguous, when their exact normalized phrase appears
in the target fields, or when an action cannot be parsed. Excluded queries
remain candidate distractors. The resulting experiment contains
\QueryCount{} queries and \CandidateCount{} candidates.

The benchmark reports Hit@$k$ for $k\in\{1,3,10\}$ and mean reciprocal rank
(MRR). Targets with non-positive scores receive zero credit. For a positive
score tied across ranks $a$ through $b$, Hit@$k$ receives the fraction of those
positions at or above the cutoff. Reciprocal-rank credit is
\begin{equation}
 \frac{1}{b-a+1}\sum_{r=a}^{b}\frac{1}{r}.
\end{equation}
Retrieval ties use a numerical tolerance of $10^{-12}$. Confidence intervals
resample whole source-video clusters, with 2,000 draws and seed 1729. The
within-video analysis restricts candidates to the query's source video.

All four terminology configurations use the same ordered query cohort and
index artifacts. New exact normalized query--target phrase matches are
recorded separately; none occurred in these runs. This check does not remove
all semantic dependence between a clip's annotation fields.

\subsection{Constructed pairwise development experiments}
The metadata-contrast dataset contains 800 steps and 2,395 comparisons selected
using parsed metadata rather than retrieval scores. The hard-negative dataset
contains 350 steps and 1,309 comparisons. It uses paraphrased requests and
alternatives selected using frozen scorer outputs. Both datasets designate
the clip supplying the request information as the winner. The hard-negative
selection is deliberately challenging and is not a random sample of user
requests.

The evaluator scores the supplied A/B clips directly. Agreement credit is one
when the designated winner receives a higher score, one-half for an exact
score tie, and zero otherwise. Intervals resample complete step clusters,
using 5,000 draws and seed 42. The October 6 run copied the earlier comparison
inputs unchanged, rather than regenerating winners after the code changes.
The two datasets are reported separately because their selection procedures
differ.

\subsection{Independent human preference evaluation}
\pending{Recruitment or annotation source, task instructions, independent A/B
judgments, disagreement handling, missing-reference counts, and the frozen
primary evaluation. No results are available yet.}

\subsection{Prospective held-out evaluation}
\pending{A new video pool and independently authored queries or judgments,
with source overlap checks and a frozen evaluation protocol. The current
250-video sample has already informed development and must not be described
as an untouched test set.}

\section{Results}
\subsection{Retrieval performance}
Lexical matching is strongest on the whole-corpus benchmark
(Table~\ref{tab:global}). Its Hit@1 is \LexicalHitOne\%, compared with
\StructuredHitOne\% for structured matching and \HybridHitOne\% for the
hybrid. The hybrid-minus-lexical Hit@1 difference is approximately $-12.6$
percentage points, with a video-clustered 95\% interval of approximately
$[-18.4,-6.9]$ points.

\begin{table}[htbp]\centering\small
\caption{Whole-corpus retrieval. Hit rates are percentages.}\label{tab:global}
\begin{tabular}{lrrrr}\toprule
Method & Hit@1 & Hit@3 & Hit@10 & MRR\\\midrule
\GlobalRows
\bottomrule\end{tabular}
\end{table}

Restricting retrieval to the same video improves all methods
(Table~\ref{tab:within}). Lexical and hybrid Hit@1 are both approximately
72.7\%. The hybrid-minus-lexical interval is approximately $[-4.9,+5.4]$
percentage points. Equal point estimates do not establish equivalence.

\begin{table}[htbp]\centering\small
\caption{Within-video retrieval. Hit rates are percentages.}\label{tab:within}
\begin{tabular}{lrrrr}\toprule
Method & Hit@1 & Hit@3 & Hit@10 & MRR\\\midrule
\WithinRows
\bottomrule\end{tabular}
\end{table}

\subsection{Terminology ablation}
Parser cleanup preserves the observed Hit@1 rates. Adding the active aliases
also leaves aggregate Hit@1 unchanged (Table~\ref{tab:ablation}). For lexical
ranking, one query moves from first to second place and another moves from
second to first. Removing substitute lists reduces observed lexical Hit@1;
the combined setting has the lowest lexical point estimate.

\begin{table}[htbp]\centering\small
\caption{Hit@1 after cleanup, with independent terminology options. Values are
percentages; hybrid weights remain fixed.}\label{tab:ablation}
\begin{tabular}{lrrr}\toprule
Configuration & Lexical & Structured & Hybrid\\\midrule
\AblationRows
\bottomrule\end{tabular}
\end{table}

\begin{table}[htbp]\centering\small
\caption{Paired lexical Hit@1 changes relative to the cleaned baseline, in
percentage points. Intervals resample source-video clusters.}\label{tab:deltas}
\begin{tabular}{lrrr}\toprule
Configuration & Change & 95\% lower & 95\% upper\\\midrule
\DeltaRows
\bottomrule\end{tabular}
\end{table}

The intervals in Table~\ref{tab:deltas} include zero. Thus the ablation does
not establish a general improvement or decline beyond this development
sample. The results support keeping the new search options separate and
optional while evaluating broader terminology coverage.

\subsection{Pairwise development agreement}
All 3,704 comparisons resolve to canonical clips. Table~\ref{tab:pairs}
reports agreement with the constructed winners. The metadata-contrast task
is nearly saturated by lexical matching. The hard-negative task separates
the methods more strongly: lexical agreement remains high while structured
agreement falls to approximately 50.8\%. The hybrid is between them.
These are target-recovery results, not observed human preference agreement.

\begin{table}[htbp]\centering\small
\caption{Agreement with constructed pairwise winners. Agreement and confidence
intervals are percentages; intervals are clustered by step.}\label{tab:pairs}
\begin{tabular}{llrrr}\toprule
Dataset & Method & Pairs & Agreement & 95\% interval\\\midrule
\PairwiseRows
\bottomrule\end{tabular}
\end{table}

\subsection{Diagnostic findings}
Table~\ref{tab:diagnostics} records overlapping diagnostic flags. Structured
ties and zero evidence are frequent, especially in the hard-negative corpus.
These counts identify cases for inspection; they do not, by themselves,
establish whether a failure was caused by missing source information, an
incorrect parse, an inadequate lexicon, or the scoring rule.

A further audit partitions the 415 comparisons in which both structured
scores are zero. In 21 comparisons, the query has no parsed action. In 58,
one candidate has no parsed actions. In the remaining 336, the query and both
candidates have parsed actions, but both selected alignments have zero action
match. Thus about 81\% of these zero-score comparisons still contain parsed
actions on both sides. This narrows the next review to the extracted actions
and their compatibility, rather than treating every zero as an absent parse.
It does not establish a synonym error: negation mismatches or genuinely
different actions can also produce zero action match.

\begin{table}[htbp]\centering\small
\caption{Diagnostic counts, reported separately by construction procedure.
Rows overlap and must not be summed as distinct errors.}\label{tab:diagnostics}
\begin{tabularx}{\linewidth}{Xrr}\toprule
Diagnostic & Metadata contrast & Hard negatives\\\midrule
\DiagnosticRows
\bottomrule\end{tabularx}
\end{table}

\subsection{Human-reviewed extraction accuracy}
\pending{Complete the 60-row extraction worksheet; report correctness of
actions, objects, and tools, along with label coverage and review limitations.
The worksheet is stratified for diagnosis, not a simple random sample for
estimating corpus-wide precision.}

\subsection{Independent preference and final held-out results}
\pending{Insert validated results from independent human judgments and the
prospective evaluation. Leave unmeasured quantities blank.}

\section{Discussion and limitations}
The current evidence favors lexical matching. Adding structure does not help
simply because the representation is more explicit: the parser must identify
the right information, and the matching rule must use it appropriately. The
large number of ties and missing-evidence cases indicates that these conditions
are often not met by the present system.

Terminology inconsistency is a concrete data issue, but it is not a complete
explanation for the performance gap. The active lexicon is small, and some
synonyms already occur in alternative fields. The ordinary clip-name benchmark
also does not specifically sample regional language or unfamiliar synonyms.
A more informative test would include independently written requests using
those variants, together with confusing non-synonyms.

Several limitations constrain interpretation. First, query names and candidate
descriptions describe the same clip, so their wording and meaning are not
independent even after direct phrase exclusions. Second, candidates with
narrative evidence are a selected subset of the full corpus. Third, parser
coverage is measured but parser precision is not. Fourth, constructed pairwise
winners and scorer-selected hard negatives do not represent a population of
human preferences. Finally, these videos and results have already been examined
during development. A new split of the same data cannot reverse that exposure.

End-state matching remains a separate, unfinished task. An intended result
such as a watertight seal is different from evidence that the seal was achieved.
Likewise, ``smooth,'' ``level,'' ``dry,'' and ``cured'' describe different
conditions. No validated end-state representation or evaluation is reported.

\section{Conclusion}
The completed experiments establish a reproducible comparison of lexical,
structured, and hybrid matching on human-written instructional video
annotations. Lexical matching is currently the strongest baseline.
Inventory cleanup improves data consistency, while the tested synonym rules
do not improve aggregate Hit@1. The next empirical requirement is independent
evaluation of parser correctness and clip preference, followed by a
prospective test of any revised method.

\pending{Revise the final conclusion after independent preference and held-out
results are available. Do not replace the present negative result with an
unmeasured improvement claim.}

\section*{Acknowledgments}
\pending{Author-approved acknowledgments and funding information.}

\bibliographystyle{plainnat}
\bibliography{references}

\appendix
\section{Reproducibility}
The verified development run is stored at
\path{project1_outputs/experiment-suite-2026-10-06/}. Its summary records the
individual configurations, paired intervals, index profile, and pairwise
diagnostics. The run manifest records hashes of inputs, Python source, and
output artifacts. The paper tables are generated from the verified JSON
summary rather than transcribed manually.

From the repository root, the full development suite can be run with:
\begin{quote}\small\ttfamily
./scripts/run\_local\_pipeline.sh experiments \\
\hspace*{1em}--output-dir project1\_outputs/new-run \\
\hspace*{1em}--development-root project1\_outputs
\end{quote}
Each run requires a new output directory. Supplying the development root
copies the existing labels unchanged; omitting it constructs new development
datasets. These two modes must be distinguished when comparing runs.

\section{Outstanding study materials}
A review packet has been prepared with the existing 60 extraction examples
and 44 diagnostic preference comparisons. The latter sample contains one
comparison per selected source step, with randomized A/B order and no model
scores or constructed winners shown in the worksheet. It includes eight
query-parse cases and twelve cases from each of three further diagnostic
groups. These groups guide case selection; they do not define the true cause
of an error. All human response fields remain blank.

Preference responses allow A, B, tie, and unjudgeable, with a blank response
treated as pending. The export tool preserves tie, unjudgeable, and pending
counts while exporting only observed A/B choices to the binary evaluator.
The sample is intended for diagnosis, so agreement on these cases must not
be presented as the final independent preference result.

A prospective source-overlap check compares database identifiers, YouTube
identifiers, and normalized URLs with the development export. The current
250-video export fails this check against itself, as expected. Passing this
check would establish only that these identifiers are disjoint; it would not
detect every reupload or prove that queries and judgments were unseen.

\pending{Final annotation protocol approval, annotator details, prospective
test data and manifest, and release or access conditions for the private data.}
\end{document}
